\chapter{Multiple Qubits}

\section{ More than the sum of its parts}

If quantum computing was nothing more than a collection of qubits processed individually, there would be little point in building a quantum computer. A single qubit can be modelled using 2x2 matrices - you can even work these out by hand! Even if we were to have a large number of these in parallel, it would still be very easy to simulate a large number of qubits. You might prefer to use a classical computer that is very efficient at evaluating linear algebra rather than pen \& paper. 

What makes qubits difficult to simulate is the interactions between them. Having qubits interact with one another allows for more complex algorithms. Thanks to entanglement, one qubit can control what happens to another. This has some interesting implications for quantum algorithms. 

\section{Reading binary}

To describe many qubits, we need to add something to Dirac notation. For instance, if we were to label our qubits alphabetically, we'd run out of letters after just 26 qubits. Thankfully, we can leverage the binary system from classical computing to describe as many qubits as we want to. Everyone is familiar with the decimal system where counting goes as 0,1,2,...,8,9 and then we repeat from 10,11,12,... . Computers, classical or quantum $^*$ count in base 2. Binary works the same way as the decimal system but instead the count repeats every 2. Numbers go as    

0,1,10,11,100,101,110,111 

Which in decimal would be 

0,1,2,3,4,5,6,7 

To read a binary number:  
  
1. Count the number of digits in the number and take this number -1 as n
2. Starting from the first digit on the left, if that digit is 1, calculate $2^n$. 
3. Moving to the next digit on the right, if that digit has a 1, calculate $2^{n-1}$. 
4. Add the result of step 3 to step 2. 
5. Repeat this process for each digit decreasing the power by 1 each time until the last digit (which should have a $2^0 = 1$ calculation). An example of this is shown below


\includegraphics[scale=0.5]{images/binary_conversion.png}
[1]

1. Here there are 8 digits, so we take $n = 8 - 1 = 7 $. 
2. For the first digit there is a 0, so we do nothing. 
3. The second digit is 1 so we do $2^{7-1} = 2^6 = 64$. 
4. Adding $0 + 64 = 64$.
5. The sum we get is

$$ 0 + 2^6 +0  + 2^4 + 2^3+ 0 + 0 + 2^0 = 89 $$

Here we needed at least 6 digits in binary to represent a number with 2 digits in decimal. The decimal notation is much more efficient. Let's say we wanted to encode the number 5 onto a quantum state. 

We know $ 5 = 2^2 + 2^0$ so we can write 5 as 101 in binary. 

As there are 3 digits, we need 3 qubits to represent this number. 

The qubit on the left is in state 1, the qubit in the middle is in state 0 and the qubit on the right is in state 1. 

$$ \ket{101} = \ket{q_2q_1q_0} $$

Here we use the convention where the first qubit ($2^0$) is on the right, with the next qubit one to the left and the last qubit the most to the right. This is what IBM uses and what we will use when programming quantum computers.

\textbf{Quantum states are read right-to-left with the first qubit next to the $\>$}

This right to left notation can also be represented as a ket as $\ket{101}_{2,1,0} = \ket{1}_2\ket{0}_1\ket{1}_0$ where the subscript index specifies which qubit is in which state. Similarly, we can have our quantum gates given these subscripts so we can say which qubits they are acting on. 

\section{ Multiple qubits, multiple gates }

If we have multiple qubits, we can apply single qubit gates on each of them at a time. This is the same as having a bunch of single qubits on their own, each having their state changed by the qubit gate. The simplest case is two qubits. Let's begin with our qubits in the normal starting state of $\ket{q_1q_0} = \ket{00}$. We can apply a Hadamard gate on the first qubit, $q_0$ to change it to the superposition state $\ket{+}$

$$ \ket{00} \rightarrow \ket{0+} $$ 

Then we might want to apply the not (X) gate to the other qubit, $q_1$ to change it to $\ket{1}$

$$ \ket{0+} \rightarrow \ket{1+} $$

As these are two separate operations, we could do them both at the same time in one step 

$$  \ket{00} \rightarrow \ket{1+} $$

\textbf{Applying single qubit gates to a collection of qubits, changes the states of the qubits the gates are applied to}

\section{ Product States }

So far we have acted on all our qubits separately, what happens to one doesn't affect the others. If we were to measure one of the qubits, it wouldn't tell us anything about the other qubits. 

Such states where we can separate the qubits out are called product states (or pure states). 

In general a product state can be written as 

$$ \ket{\Psi} = \ket{\psi_0}\ket{\psi_1}...\ket{\psi_{n-1}}$$

Let's say we measure the first qubit (qubit 0) and get a 0. The remaining qubits wouldn't be affected, so our new state would be 

$$ \ket{\Psi} = \ket{0}\ket{\psi_1}...\ket{\psi_{n-1}}$$

\section{ The Tensor Product }

Mathematically there is a very satisfying way of representing this with linear algebra. The tensor product allows us to do this. At this point some subscripts can be useful for labelling our qubits. The first (or zeroth) qubit can be denoted by a small $a$ on the right of the ket, the second (or first) qubit will have a little $b$ on the right. If qubits $a$ and $b$ are not entangled, the system of $ab$ can be described as a tensor product. 
$$ \ket{\psi_a}_a\ket{\psi_b}_b = \ket{\psi_a}_a \otimes \ket{\psi_b}_b $$
Where the symbol $\otimes$ represents the tensor product operation. 

Expanding our kets into column vectors

$$ \ket{\psi_a}_a = \begin{bmatrix} a_0 \\ a_1 \end{bmatrix} ; \ket{\psi_b}_b = \begin{bmatrix} b_0 \\ b_1 \end{bmatrix} $$

$$
 \ket{\psi_a}_a \otimes \ket{\psi_b}_b = \begin{bmatrix} a_0\ket{\psi_b}_b \\ a_1\ket{\psi_b}_b \end{bmatrix}
$$

This is a new vector with $\ket{\psi_b}_b$ repeated twice, each time multiplied by a scalar $b_0$ or $b_1$. The final vector has 4 elements and looks like: 
$$
 \ket{\psi_a}_a \otimes \ket{\psi_b}_b = \begin{bmatrix} a_0 b_0 \\ a_0 b_1 \\ a_1b_0 \\ a_1 b_1 \end{bmatrix}
$$

So you can think of your new state as a 4-level quantum system. Every time we add another qubit, we tensor product it in, so the size of the vector doubles. This is where the exponential scaling of the state-space comes from. The same applies for multiple quantum gates in parallel. Our matrix increases by a factor of 2 in both the number of rows or columns. This is even better scaling!

Keeping with the number 5, the state vector for this number 

$$ \ket{5} = \ket{101} = \begin{bmatrix} 0 \\ 0\\ 0\\ 0\\ 0\\ 1\\0\\0 \end{bmatrix} $$

For each qubit we add, the number of rows in this column vector doubles. This is what makes it difficult to simulate quantum algorithms on classical computers. 

\section{ Entanglement } 

You might have noticed that it's really easy to apply single qubit gates to a collection of qubits. In the last chapter, we saw that all single qubit gates can be represented by $ 2 \times 2$ matrices. And each qubit can be represented by a column vector with 2 rows. So if we had a large number of qubits, we could still simulate a large number of $2 \times 2$ matrices on a classical computer. If it is easy to simulate a collection of isolated qubits, why do we bother making quantum computers? 

Alone superposition has quite limited applications. Applying single qubit operations can be used for sampling distributions, but this limits quantum computers to being random number generators. Where superposition can become more useful is when entanglement enters the picture. 

\subsection{ If statements with spicy sauce }

In general programming, if statements allow you to include logic in a program. If a statement is true, the algorithm will do one thing, otherwise, if the statement is false, it will do something else. The computer first has to check the variables value before executing the next step of the algorithm. We can do this on a quantum computer by measuring our qubits, collapsing the superposition and then performing our next quantum gates accordingly. For example, \href{https://qiskit.org/textbook/ch-algorithms/teleportation.html}{in quantum teleportation this is done}. 

Back in section\href{casino} there was a game where if a coin landed up we got +1 dollar, otherwise (if it were tails up) we got -1 dollars. This works as a classical if statement, we add the value after the coin is landed and observed. What makes quantum operators different is that we can apply them to the superposition. What this means is we can act on both the $\ket{H}$ and $\ket{T}$ amplitudes at once. 

All the way back in Chapter 1, entanglement was introduced 
 
 In general English, we use the word entangled to suggest that a process has mixed things up in such a way that they are hard to separate. This definition works nicely for quantum states where we cannot describe one qubit without describing another. What this means, is that when two qubits are entangled, the superposition of one qubit defines the superposition of another. Mathematically, we can't treat our qubits as a product state of separate 2x1 vectors any more. 

\subsection{ Can the qubit not? }

Entanglement can be created by gates that act on multiple qubits together to create an entangled state. The most common way this is done is by using a 2 qubit gate called the controlled-not, usually shown as CNOT or CX. The CNOT acts on two qubits, the control qubit and the target qubit. If the control qubit is in the state $\ket{0}$, it will do nothing, if the control qubit is in the state $\ket{1}$ it will perform the not (X) gate on the target qubit. In both cases, the state of the control qubit is unchanged. We can represent this with a table

\begin{center}
\begin{tabular}{ | c | c | c | c | c |}

 Control & Target & $\rightarrow{}$ & Control & Target \\ 
  
 0 & 0 &  & 0 & 0 \\ 
 
 0 & 1 &  & 0 & 1 \\
 
 1 & 0 &  & 1 & 1 \\
 
 1 & 1 &  & 1 & 0 \\

  
\end{tabular}
\end{center}

If we have two qubits $\ket{psi_b}_b\ket{psi_a}_a$ where qubit 0 is in the state $\ket{psi_a}_a$ and qubit 1 is in the state $\ket{psi_b}_b$, there are two CNOT's we can have. One where qubit 0 is the target and qubit 1 is control or the other way around, qubit 1 is the control and qubit 0 is the target. Each of these has a different 4x4 matrix

The first has the control on the 
$$
CNOT_{0, 1} = \begin{bmatrix} 1 & 0 & 0 & 0 \\
0 & 0 & 0 & 1 \\
0 & 0 & 1 & 0 \\
0 & 1 & 0 & 0 \\
\end{bmatrix}
$$

$$
CNOT_{1, 0} = \begin{bmatrix} 1 & 0 & 0 & 0 \\
0 & 0 & 0 & 0 \\
0 & 0 & 0 & 1 \\
0 & 0 & 1 & 0 \\
\end{bmatrix}
$$

\subsection{ Ring the Bell (states)}

There are many ways we can entangle two qubits, and some qubits can be more entangled than others. The most common entangled states are known as the Bell states, and can be written like

$$ \ket{\Phi^+} = \frac{1}{2} \left( \ket{00} + \ket{11} \right) $$
$$ \ket{\Phi^-} = \frac{1}{2} \left( \ket{00} - \ket{11} \right) $$
$$ \ket{\Psi^+} = \frac{1}{2} \left( \ket{01} + \ket{10} \right) $$
$$ \ket{\Psi^-} = \frac{1}{2} \left( \ket{01} - \ket{01} \right) $$

For each of these states, we can't separate out our two qubits as a product state. 


The following circuit generates the Bell states

\includegraphics{images/Bell_State_circuit.png}


\subsection{General controlled unitaries }

Instead of having a not-gate being applied, we can have any single quantum gate. The logic follows from the controlled not. For the control qubit in the state $ket{0}$ no action is carried out, whilst the control qubit in the state $\ket{1}$ will implement the unitary on the target. 

\subsection{Multiple-control qubit gates }

We can have controlled quantum gates where there are more than one qubit as control. In this case, the target qubit is only affected if all the control qubits are 1. 

The Toffoli gate is the same as the controlled-not gate but with two qubits needing to be $\ket{1}$. Only if both the control qubits are in the state $\ket{1}$ for the $X$ to apply on the target qubit. Otherwise the Toffoli gate acts as identity. 

\section{Chapter Exercises}

\begin{enumerate}
    \item What would the following numbers be in decimal?

a. 1111

b. 1001

c. 101101

d. 0110


    \item What would the following numbers be in binary? 

a. 1 

b. 5

c. 6 

    \item For n qubits, we need a vector with $2^n$ rows to describe the state. What shape of matrix (rows, columns) do we need to describe the operator (i.e. multi-qubit gate) that acts on this vector? 
\end{enumerate}

\section{Chapter Summary}

\begin{itemize}
    \item Quantum computers work in binary
    \item Many qubits have quantum states represented as a binary number
    \item Quantum states are read right-to-left with the first qubit at the end
    \item We can apply single qubit gates individually across multiple qubits 
    \item When qubits are independent of one another, they are in a product state
    \item When one qubit influences the state of another qubit they are entangled 
    \item We can use entanglement to apply logic gates
\end{itemize}

$^*$ There have been proposals for using a base other than 2 for quantum computing. This has some advantages but the difficulty in constructing algorithms combined with the difficulties in error correction such a device has prevented any significant machine from being built.

\section{Additional reading}

[1] L.Vivah,\href{https://medium.com/@LindaVivah/learn-how-to-read-binary-in-5-minutes-dac1feb991e}{Learn How to Read Binary in 5 minutes}, 2018
